\documentclass[12pt]{article}
\usepackage{amsmath,amssymb}
\usepackage{geometry}
\geometry{a4paper, margin=2.5cm}

\begin{document}
\section*{Fourier Neural Operator}

\subsection*{Постановка задачи}

Рассматривается задача аппроксимации оператора решения уравнений в частных производных. Пусть
\[
\mathcal{G} : \mathcal{A} \to \mathcal{U}
\]
— нелинейный оператор, отображающий входную функцию
\[
a(x) \in \mathcal{A}
\]
(например, параметры среды, начальные или граничные условия)
в решение
\[
u(x) \in \mathcal{U},
\]
удовлетворяющее некоторому уравнению в частных производных.

Целью является построение нейросетевой аппроксимации оператора $\mathcal{G}$, устойчивой к изменению дискретизации и способной работать на произвольных сетках.

\subsection*{Идея операторного обучения}

В отличие от стандартных нейросетей, аппроксимирующих функции конечномерных аргументов, Fourier Neural Operator (FNO) аппроксимирует отображения между функциональными пространствами. Таким образом, обучается не конкретное решение, а сам оператор, связывающий входные и выходные поля.

\subsection*{Спектральное представление}

Пусть $u(x)$ — поле, определённое в области $\Omega \subset \mathbb{R}^d$. Рассмотрим его преобразование Фурье:
\[
\hat{u}(k) = \mathcal{F}[u](k),
\]
где $k$ — волновой вектор. В спектральном пространстве дифференциальные операторы принимают алгебраическую форму, что делает представление глобальных зависимостей более эффективным.

\subsection*{Слой Fourier Neural Operator}

Один слой FNO определяется следующим образом:
\[
u^{(l+1)}(x)
=
\sigma\!\left(
\mathcal{F}^{-1}
\left(
W_k \, \mathcal{F}\left(u^{(l)}\right)(k)
\right)(x)
+
B \, u^{(l)}(x)
\right),
\]
где:
\begin{itemize}
\item $u^{(l)}(x)$ — поле на $l$-м слое,
\item $\mathcal{F}$ и $\mathcal{F}^{-1}$ — прямое и обратное преобразования Фурье,
\item $W_k$ — обучаемые матрицы весов в спектральном пространстве, определённые для ограниченного числа низкочастотных мод,
\item $B$ — линейный оператор в физическом пространстве,
\item $\sigma(\cdot)$ — нелинальная функция активации.
\end{itemize}

Таким образом, слой FNO сочетает глобальное спектральное преобразование с локальной линейной коррекцией.

\subsection*{Архитектура сети}

Полная архитектура Fourier Neural Operator состоит из следующих этапов:
\begin{enumerate}
\item \textbf{Lifting layer} — проекция входного поля в пространство большей размерности:
\[
u^{(0)}(x) = P(a(x)).
\]
\item \textbf{Последовательность FNO-слоёв}:
\[
u^{(0)} \to u^{(1)} \to \dots \to u^{(L)}.
\]
\item \textbf{Projection layer} — отображение в пространство решений:
\[
u_{\text{out}}(x) = Q(u^{(L)}(x)).
\]
\end{enumerate}

\subsection*{Обучение}

Обучение сети осуществляется на выборке пар
\[
\{ a^{(i)}(x), u^{(i)}(x) \}_{i=1}^{N},
\]
где $u^{(i)}(x)$ — решения, полученные с помощью численных методов (например, FEM или FDM).

Функционал потерь обычно выбирается в виде нормы $L^2$:
\[
\mathcal{L}
=
\frac{1}{N}
\sum_{i=1}^{N}
\left\|
u_{\text{pred}}^{(i)} - u_{\text{true}}^{(i)}
\right\|_{L^2(\Omega)}^2.
\]

\subsection*{Свойства метода}

Fourier Neural Operator обладает следующими характеристиками:
\begin{itemize}
\item аппроксимация операторов, а не отдельных решений;
\item независимость от конкретной пространственной сетки;
\item высокая эффективность для волновых и диффузионных задач;
\item высокая скорость инференса после обучения.
\end{itemize}

В то же время метод требует наличия обучающей выборки и не обеспечивает строгого выполнения физических законов без дополнительного введения физически-информированных ограничений.

\begin{table}[h]
\centering
\caption{Сравнение численных и нейросетевых методов решения PDE}
\begin{tabular}{|l|c|c|c|}
\hline
\textbf{Критерий} & \textbf{FEM / FDM} & \textbf{PINN} & \textbf{FNO} \\
\hline
Тип аппроксимации & Решение & Решение & Оператор \\
\hline
Использование физики & Полное & Полное (в loss) & Косвенное \\
\hline
Требование к данным & Не требуется & Минимальное & Большое \\
\hline
Работа с волнами & Высокая точность & Ограниченная & Высокая \\
\hline
Масштабируемость (3D) & Ограниченная & Низкая & Высокая \\
\hline
Скорость инференса & Низкая & Низкая & Очень высокая \\
\hline
Подходит для обратных задач & Ограниченно & Да & Ограниченно \\
\hline
Зависимость от сетки & Да & Нет & Нет \\
\hline
\end{tabular}
\end{table}
